%% kapa: exam=2 round=29 year=2018 period=3 subject=compensation_law subject_ko=감정평가 및 보상법규 seq=2 number=2 points=30
%% source: Q-Net 감정평가사 기출문제 (cst00309) artlSeq=5203398 file=raw/1008/5203398/3교시 감정평가 및 보상법규.pdf
%% conversion: kordoc 4.24.51 -> tools/gen2.py (auto); review=proofread-gpt-terra; flags=none
\begin{longquestion}{【문제 2】}
甲은 2014. 3.경 감정평가사 자격을 취득한 후, 2015. 9. 2.부터 2017. 8. 3.까지 ‘乙 감정평가법인’의 소속 감정평가사였다. 또한 甲은 2015. 7. 7.부터 2017. 4. 30.까지 ‘수산업협동조합 중앙회(이하 ‘수협’이라 함)’에서 상근계약직으로 근무하였다. 관할 행정청인 국토교통부장관 A는 甲이 위와 같이 수협에 근무하면서 일정기간 동안 동시에 乙 감정평가법인에 등록하여 소속을 유지하는 방법으로 감정평가사 자격증을 대여하거나 부당하게 행사했다고 봄이 상당하여, 「감정평가 및 감정평가사에 관한 법률」(이하 ‘감정평가법’이라 함) 제27조가 규정하는 명의대여 등의 금지 또는 자격증 부당행사 금지에 위반하였다는 것을 이유로 징계처분을 내리고자 한다. 다음 물음에 답하시오. (30점) \par
1) 국토교통부장관 A가 甲에 대하여 위와 같은 사유로 감정평가법령상의 징계를 하고자 하는 경우, 징계절차에 관하여 설명하시오. (20점) \par
2) 위 징계절차를 거쳐 국토교통부장관 A는 甲에 대하여 3개월간의 업무정지 징계처분을 하였고, 甲은 당해 처분이 위법하다고 보고 관할법원에 취소소송을 제기하였다. 이 취소소송의 계속 중 국토교통부장관 A는 당해 징계처분의 사유로 감정평가법 제27조의 위반사유 이외에, 징계처분 당시 甲이 국토교통부장관에게 등록을 하지 아니하고 감정평가업무를 수행하였다는 동법 제17조의 위반사유를 추가하는 것이 허용되는가? (10점) \par
\begin{framedtext}
참조 조문 \par
<감정평가 및 감정평가사에 관한 법률> 
제17조(등록 및 갱신등록) ① 
제11조에 따른 감정평가사 자격이 있는 사람이 
제10조에 따른 업무를 하려는 경우에는 대통령령으로 정하는 기간 이상의 실무수습을 마치고 국토교통부장관에게 등록하여야 한다. 
②\textasciitilde{}
④ 생략 
제27조(명의대여 등의 금지) 감정평가사 또는 감정평가업자는 다른 사람에게 자기의 성명 또는 상호를 사용하여 
제10조에 따른 업무를 수행하게 하거나 자격증ㆍ등록증 또는 인가증을 양도ㆍ대여하거나 이를 부당하게 행사하여서는 아니 된다.
\end{framedtext}

\end{longquestion}
